% GENERADO por benchmarks/np2_sequence/tesis/etapa3_tablas.jl; no editar a mano.
% Procedencia:
%   carbon_rohf_results_3P_R30.0.jld2 (sha256 e8c5884b)
%   CI carbon_rohf_results_3P_R30.0.jld2 m=40 (sha256 e8c5884b, commit 6dfb84b)
%   silicon_rohf_results_3P_R30.0.jld2 (sha256 09adeb32)
%   CI silicon_rohf_results_3P_R30.0.jld2 m=20 (sha256 09adeb32, commit 2479092)
%   germanium_rohf_results_3P_R30.0.jld2 (sha256 7567c491)
%   CI germanium_rohf_results_3P_R30.0.jld2 m=20 (sha256 7567c491, commit 2479092)
%   tin_rohf_results_3P_R30.0.jld2 (sha256 7db46b8f)
%   CI tin_rohf_results_3P_R30.0.jld2 m=20 (sha256 7db46b8f, commit 2479092)
%   zeta en cm^-1. zeta_desnudo = (alpha^2/2) Z <r^-3>. zeta_NIST: la zeta unica que reproduce 3P_2, 1D_2 y 1S_0 del NIST en Breit-Pauli de p^2 sin CI. R = (3P_2 - 3P_1)/(3P_1 - 3P_0); una zeta de un cuerpo en LS puro da R = 2.
\begin{tabular}{lrccccccc}
    \toprule
    \textbf{Elemento} & $Z$ & $(\alpha Z)^2$ & $\zeta_{\text{desnudo}}$ & $\zeta_{\text{calc}}$ & $\zeta_{\text{NIST}}$ & $\zeta_{\text{calc}}/\zeta_{\text{NIST}}$ & $R_{\text{NIST}}$ & $R_{\text{modelo}}$ \\
    \midrule
    Carbono & 6 & 0.0019 & 59.3 & 42.2 & 28.9 & 1.460 & 1.644 & 1.983 \\
    Silicio & 14 & 0.0104 & 167.5 & 141.0 & 148.1 & 0.952 & 1.894 & 1.911 \\
    Germanio & 32 & 0.0545 & 896.3 & 823.7 & 920.6 & 0.895 & 1.531 & 1.563 \\
    Estaño & 50 & 0.1331 & 1992.5 & 1885.9 & 2241.1 & 0.842 & 1.026 & 1.137 \\
    \bottomrule
\end{tabular}
